% Graphon-MARL robotics experiment writeup.
% Include this file from the graphon-marl repository root, or update the
% figure paths if it is included from another directory.

\paragraph{Scaling experiment ($n=1{,}000$).}
Agents are arranged on a $25\times 40$ grid, yielding $n=1{,}000$ agents.
We sweep $\kappa \in \{1,3,5,10,15,20,25,30,60,100\}$.
Each configuration is trained for $T=250$ value-iteration steps with $M=50$
Monte Carlo samples per Bellman update, and evaluated over $5$ independent
training seeds with $20$ rollouts per seed ($100$ total returns per $\kappa$).

\paragraph{Noise robustness experiment ($n=100$).}
Agents are arranged on a $10\times 10$ grid, yielding $n=100$ agents. We sweep
\[
\kappa \in \{1,5,10,15,20,25\},
\]
and evaluate three graphon conditions:
(i) \textbf{clean}, using the true graphon weights;
(ii) \textbf{Gaussian noise}, where
$\varepsilon_{ij}\sim\mathcal{N}(0,\sigma^2)$ is added to each entry of the
row-normalized weight matrix; and
(iii) \textbf{subgaussian noise}, where
$\varepsilon_{ij}\sim \mathrm{Uniform}(-\sigma,\sigma)$, which is subgaussian
with parameter $\sigma/\sqrt{3}$.
In the noisy settings, perturbed weights are clipped to $[0,\infty)$ and
renormalized. We use $\sigma=0.05$ and apply noise only at evaluation time,
while training is always performed on the clean graphon. Each condition uses
$5$ training seeds with $50$ rollouts per seed ($250$ total returns per
$\kappa$).

Figure~\ref{fig:robotics_scaling_noise} reports the aggregate return curves.
The sampled-neighborhood perception maps in
Figure~\ref{fig:robotics_perception_heatmap} visualize how increasing
$\kappa$ changes the empirical local neighborhood observed by a representative
agent.

\paragraph{Interpreting the population-size comparison.}
The comparable clean returns for $n=100$ and $n=1{,}000$ at shared values of
$\kappa$ are expected in this setup. The reported objective is a discounted
\emph{per-agent average} return, not total system reward. The graphon weights
are row-normalized, so each agent observes a local probability distribution
over neighbor states; increasing $n$ refines the empirical approximation of the
same radial graphon rather than changing the scale of an individual agent's
reward. The GMFS offline table is indexed by the local state, action, and
$\kappa$-sampled neighborhood histogram, so $n$ enters online evaluation
mainly through finite-population and boundary effects. Thus larger $n$ should
mostly reduce discretization noise and need not produce a large shift in mean
per-agent return for fixed $\kappa$.

\begin{figure}[t]
    \centering
    \begin{subfigure}[t]{0.55\linewidth}
        \centering
        \includegraphics[width=\linewidth]{results_robotics_n1000/plots/discounted_reward_vs_k.pdf}
        \caption{$n=1{,}000$: return versus subsample size $\kappa$.}
        \label{fig:reward_vs_k_n1000}
    \end{subfigure}
    \hfill
    \begin{subfigure}[t]{0.44\linewidth}
        \centering
        \includegraphics[width=\linewidth]{plots_n100_noise/n100_noise_comparison.pdf}
        \caption{$n=100$: clean and noisy graphon evaluation.}
        \label{fig:n100_noise_comparison}
    \end{subfigure}
    \caption{\textbf{Robotics coordination results for scaling and noise robustness.}
    Error bars show standard error across evaluated returns.}
    \label{fig:robotics_scaling_noise}
\end{figure}

\begin{figure}[t]
    \centering
    \includegraphics[width=0.78\linewidth]{results_robotics_n1000/plots/perception_heatmap_compare_first_10.pdf}
    \caption{\textbf{Perception heatmaps for graphon subsampling ($n=1{,}000$).}
    For a representative focal agent on the $25\times 40$ grid, each panel
    aggregates sampled neighbors over the first ten rollout steps for a
    different subsample size $\kappa$. Larger $\kappa$ gives denser empirical
    coverage of the same row-normalized radial graphon neighborhood, while the
    sampled mass remains localized around the focal agent's interaction
    region.}
    \label{fig:robotics_perception_heatmap}
\end{figure}
